%%%%%%%%%%%%%%%%%%%%%%%%%%%%%%%%%%%%%%%%%%%%%%%%%%%%%%%%%%%%%%%%%%%%%%%%%%%%%%%%% 
%%%%%%%%%%%%%%%%%%%%%%%%%%%%%%%%%%%%%%%%%%%%%%%%%%%%%%%%%%%%%%%%%%%%%%%%%%%%%%%%% 

\paragraph{Problem 1(d) answer} \GradescopeComment{This should be on page 5 (middle) of your PDF.}

\begin{theorem}
  For any instance $I=\{J_1, J_2, \ldots, J_n\}$ of \prob{Puncturing Intervals},
  \textsf{rpuncture}${(I)}$ returns a correct solution for $I$.
\end{theorem}

% Prove the theorem above for your rpuncture algorithm from Part (c).
% As much as possible, your proof should mimic the structure
% of the proof of Theorem 5.2 (for rselect) in Lecture Note 5.

\begin{blue}
  
  \begin{longFormProof}

    \step The proof is by induction on $n$.

    \step\label{step: base case}
    For the base case $n=0$, \textsf{rselect} returns the empty set, so is correct in this case.

    \smallskip 

    \begin{block}\label{block: inductive step}

      \step\label{step: inductive assumption}
      Consider executing $\textsf{rselect}(I)$ for any $I$ with $n\ge 1$.
      Assume \textsf{rselect} is correct on smaller instances.

      % [review 2026-09-27] fix: J_1, D, and R are computed in Lines 2, 3, and 4 of the pseudo-code (was "Lines 2 and 3")
      \step Let $J_1$, $D$, and $R=\textsf{rselect}(D)$ be as computed by Lines 2--4 of the algorithm.

      \step Note that $|D|<|I|$,
      so by the assumption in Step~\ref{step: inductive assumption} $R$ is a correct solution for $D$.

      \step So, by Lemma 2, the solution $\{J_1\} \cup R$ returned by $\textsf{rselect}(I)$ is correct for $I$.

    \end{block}

    \smallskip 

    \step By Block~\ref{block: inductive step},
    for any instance $I$ with $n\ge 1$,
    if \textsf{rselect} is correct on smaller instances,
    then $\textsf{rselect} (I)$ is correct.

    \step By this, Step~\ref{step: base case}, and induction, \textsf{rselect} is correct for all instances.
  \end{longFormProof}

\end{blue}

\vfill


\paragraph{Problem \arabic{problem} collaboration and citations}~\GradescopeComment{This should be on page 5 (bottom) of your PDF.}


I collaborated on this problem with the following people:
\begin{blue}
  \begin{itemize}
  \item \REPLACEME                    % replace with "none" if none
  \end{itemize}
\end{blue}

My answers use ideas from the following sources (other than the lecture notes and textbooks): 
\begin{blue}
  \begin{itemize}
  \item \REPLACEME                    % replace with "none" if none
  \end{itemize}
\end{blue}

%%% Local Variables:
%%% mode: latex
%%% TeX-master: "homework_5"
%%% End:
